% Source: https://artofproblemsolving.com/community/c3h1257865_what_is_the_value_of_infinitely_nested_0s
% Stable identity: aops-topic-1257865-q01
The equation $x=\sqrt{x}$ only tells us which numbers \emph{could} be the limit; to decide, look at the sequence. Here $a_1=0$, and then $a_2=\sqrt0=0$, and every term is $0$. So the expression equals $\boxed{0}$; the root $x=1$ is extraneous.

For $b_n$: $b_1=\sqrt6<3$, and if $b_n<3$ then $b_{n+1}=\sqrt{6+b_n}<\sqrt9=3$; also $b_{n+1}^2-b_n^2=6+b_n-b_n^2=(3-b_n)(2+b_n)>0$, so the sequence increases. An increasing sequence bounded above converges, and its limit $L$ satisfies $L=\sqrt{6+L}$, $L\ge0$, so $L=3$. This is the justification the textbook asks for: the equation identifies the limit only after the sequence is known to converge, and among the roots one must pick the one the sequence actually approaches.
